\section{An exact statistical frontier at all cuts}
\label{sec:v22-revelation}
The preceding realizations are uniform upper bounds. A family with an
exact answer illustrates what an all-algorithm statistical lower bound
requires. Unlike exact computation of a prescribed word function, the
objective here is the best decision risk among all rules, including
stochastic rules, under a fixed preparation schedule and full profile.

Let $\Theta=\{1,\ldots,M\}$ and fix a known prior
$\pi_1\ge\cdots\ge\pi_M\ge0$, $\sum_j\pi_j=1$. At time $t$, an
independent indicator $E_t$ has known mean $q_t\in[0,1]$. The report is
$Y_t=\vartheta$ if $E_t=1$, and $Y_t=*$ otherwise. All indicators are
independent of the parameter and the machine's private randomization.
The unknown parameter $\vartheta$ is fixed during the $n$ observations.
The machine starts in a parameter-independent state, has at most $K_t$
labels after report $t$, and finally guesses $\vartheta$. Its initial
and final randomization is part of the same finite causal machine.
The objective is Bayesian probability of a correct answer under $\pi$.
Parameter-dependent loss is specified here and is not identified with
the one-row physical signed-score problem.

Write $P_k=\sum_{j=1}^k\pi_j$ and put
\begin{equation}\label{eq:v22-revelation-weights}
 k_t=\min\{M,K_t,K_{t+1},\ldots,K_n\},\quad
 a=\prod_{s=1}^n(1-q_s),\quad
 w_t=q_t\prod_{s=t+1}^n(1-q_s).
\end{equation}
Thus $a+\sum_t w_t=1$ and $k_t$ is nondecreasing in $t$.

\begin{theorem}[Exact all-cut revelation frontier]
\label{thm:v22-revelation}
Over all stochastic causal machines satisfying the specified profile,
the optimal Bayesian success probability is
\begin{equation}\label{eq:v22-exact-frontier}
 \mathsf S_{\boldsymbol q,\boldsymbol K}(\pi)
       =a\pi_1+\sum_{t=1}^n w_tP_{k_t}.
\end{equation}
A single deterministic machine attains this value simultaneously for
all $\boldsymbol q\in[0,1]^n$. The upper bound therefore applies to
all approximate statistical decision rules, not just to machines
implementing a chosen selector exactly.
\end{theorem}
\begin{proof}
Let $L$ be the last time with $E_L=1$, with $L=0$ when no revelation
occurs. Then $\Pp(L=0)=a$ and $\Pp(L=t)=w_t$. These events, and every
particular pattern of indicators, are independent of $\vartheta$.

First consider any stochastic machine. Realize all of its transition
and output kernels by independent random variables indexed by time,
state and input, and also realize its initial randomization. This is
possible on a product space independent of $\vartheta$ and the
indicators. Conditional on the entire random table the machine is
deterministic. This conditioning is an analysis device, not a free
persistent random seed in an implementation.

Fix that table and an indicator pattern with last revelation $t>0$.
The map from $\vartheta$ to the terminal guess factors through the
state at every time $s\ge t$: all reports after $t$ are the same
symbol $*$, so the remaining transitions and output are fixed maps.
The set of possible terminal guesses therefore has cardinality at
most $\min_{s\ge t}K_s$, and at most $M$ guesses can be correct.
Its prior success probability is bounded by the mass of the $k_t$
largest prior atoms, namely $P_{k_t}$. Earlier revelations and any
parameter information encoded before $t$ do not evade this future
bottleneck. On the no-revelation pattern, the output is independent of
$\vartheta$ and its success probability is at most $\pi_1$.
Integrate these bounds over the random table and indicator patterns to
obtain the upper bound in \eqref{eq:v22-exact-frontier}. In particular
no extremality claim about a nonconvex minimax set has been used.

For attainment, use the labels $\{1,\ldots,k_t\}$ at time $t$, start
with label one, and use the following deterministic updates. On $*$,
hold the label. On a revelation $j\le k_t$, replace it by $j$; on a
revelation $j>k_t$, replace it by one. The holding update is legal
because $k_{t-1}\le k_t$. At termination output the label. Conditional
on $L=t>0$, exactly the parameters $1,\ldots,k_t$ are correct; on
$L=0$, exactly parameter one is correct. This proves equality. The
updates use only the prescribed profile and prior ordering, not the
probabilities $q_t$, establishing simultaneous attainment.
\end{proof}

\begin{corollary}[Sharp preparation--peak tradeoff]
\label{cor:v22-erasure-tradeoff}
For the uniform prior, a constant reveal probability $0<q<1$, and
constant width $1\le K\le M$, write $a_n=(1-q)^n$. Then
\begin{equation}\label{eq:v22-uniform-frontier}
 \mathsf S_{n,K}=\frac{a_n+K(1-a_n)}M.
\end{equation}
For $n\ge1$ and target success $s\in[1/M,1]$, the least feasible
width is
\begin{equation}\label{eq:v22-width-inverse}
 K_{n,s}=\max\left\{1,
       \left\lceil\frac{Ms-a_n}{1-a_n}\right\rceil\right\},
\end{equation}
provided this integer is at most $M$; otherwise the target is
infeasible at that preparation count. If $s>1/M$ and $K>Ms$, the
least preparation count for width $K$ is
\begin{equation}\label{eq:v22-samples-inverse}
 n_{K,s}=\left\lceil
 \frac{\log((K-1)/(K-Ms))}{-\log(1-q)}\right\rceil.
\end{equation}
If $K\le Ms$ and $s>1/M$, no finite $n$ attains the target, except
that equality can be attained when revelations are certain, a case
excluded by $q<1$. Width one gives success exactly $1/M$ for all $n$.
\end{corollary}
\begin{proof}
Here $P_k=k/M$ and $k_t=K$. Since $\sum_tw_t=1-a_n$,
Theorem~\ref{thm:v22-revelation} gives \eqref{eq:v22-uniform-frontier}.
Solving $a_n+K(1-a_n)\ge Ms$ first for $K$ and then for $n$ proves the
formulas. For $K>1$, the limiting value $K/M$ is approached strictly
from below when $0<q<1$, which proves the boundary assertion.
\end{proof}

\begin{corollary}[Exact cost of a changed checkpoint]
\label{cor:v22-cut-sensitivity}
Replacing a profile $\boldsymbol K$ by a tighter profile
$\widetilde{\boldsymbol K}$ loses exactly
\[
 \sum_{t=1}^n w_t(P_{k_t}-P_{\widetilde k_t}).
\]
Inserting a charged holding cut of width $H$ immediately after time
$s$ replaces $k_t$ by $\min(k_t,H)$ for $t\le s$, and leaves the
later terms unchanged. In particular a single narrow early cut does
not constrain information acquired after that cut.
\end{corollary}
\begin{proof}
Subtract \eqref{eq:v22-exact-frontier} for the two profiles. A holding
cut is a further bottleneck exactly for revelations that precede it.
The same deterministic construction realizes the refined schedule.
\end{proof}

For example, with $M=4$, $q_1=q_2=q_3=1/2$, and
$(K_1,K_2,K_3)=(4,1,4)$, the effective profile is $(1,1,4)$ and the
optimal uniform-prior success is $5/8$. Four labels at the terminal
cut alone would instead give $29/32$. These are risks of complete
experiments, not counts of residual classes of one word function.
For a nonuniform prior the same theorem weights every surviving label
by its actual mass. The nested optimal label sets are what makes all
cut bounds jointly attainable. For a general loss or observation
channel such nesting need not hold; compatible continuation updates,
not independently optimized dictionaries, remain necessary.

The exact formula is Bayesian. Its upper bound also obstructs any
uniform worst-case guarantee exceeding that Bayesian value, but the
attaining prior-ordered machine is not asserted to be minimax. In
particular an uncharged persistent random permutation cannot be used
to turn its Bayesian success into a minimax statement. The prior,
observation channel and all register cuts have been declared explicitly.
